A l'espectroscòpia de fotoemissió ultraviolada (UV), il·luminem les mostres amb un feix de radiació UV i analitzem l'energia dels electrons emesos.

\begin{apartat}{1}
Hem il·luminat una mostra amb radiació de longitud d'ona $\lambda = 23,7\ \mathrm{nm}$ i els fotoelectrons analitzats tenen una energia cinètica màxima de 47,7~eV. Calculeu la funció de treball del material analitzat en J i en eV.
\end{apartat}

\begin{apartat}{1}
Determineu el llindar de longitud d'ona per a aquest material. Com canviaria aquest llindar de longitud d'ona si es dupliqués la potència del feix de radiació UV?
\end{apartat}

\dades{%
  $h = 6,63\times 10^{-34}\ \mathrm{J\,s}$\\
  $1\ \mathrm{eV} = 1,60\times 10^{-19}\ \mathrm{J}$\\
  $c = 3,00\times 10^{8}\ \mathrm{m\,s^{-1}}$}
